% !TeX root = ./rpz.tex
\chapter{{\fontsize{14pt}{21pt}\selectfont\MakeUppercase{Приложение Б}}}
\label{cha:appendix_code}

\counterwithin{listing}{chapter}
\setcounter{listing}{0}

% ──────────────────────────────────────────────────────────────────────────────
% Листинг 1: алгоритм оптимизации маршрута (Go)
% ──────────────────────────────────────────────────────────────────────────────

\begin{codelisting}{Алгоритм ближайшего соседа с временными окнами (nearest\_neighbor.go)}{app:lst_optimizer}
  \begin{minted}{go}
package optimizer

import (
	"context"
	"math"
)

// NearestNeighborTW реализует эвристику ближайшего соседа для задачи
// коммивояжёра с временными окнами (NNH-TW).
//
// Алгоритм:
//  1. Выбрать начальный узел: с наиболее ранним открытием временного окна
//     или узел 0, если временные окна не заданы.
//  2. Отметить его посещённым; установить текущее время = startTime + длительность.
//  3. Повторять, пока не посещены все узлы:
//     a. Проход 1: найти ближайший допустимый узел (укладывается в окно).
//     b. Проход 2: если таких нет — взять ближайший без учёта окна.
//     c. Вычислить время ожидания, если прибываем до открытия окна.
//     d. Сдвинуть текущее время: прибытие + ожидание + длительность обслуживания.
//  4. Вернуть упорядоченные индексы узлов и агрегированную статистику.
//
// Временная сложность: O(n^2).
type NearestNeighborTW struct{}

func NewNearestNeighborTW() *NearestNeighborTW { return &NearestNeighborTW{} }

func (a *NearestNeighborTW) Name() string { return "nearest-neighbor-tw" }

func (a *NearestNeighborTW) Optimize(
	_ context.Context, g *Graph, startTimeMins int,
) (Result, error) {
	n := len(g.Nodes)
	if n == 0 {
		return Result{}, nil
	}

	visited := make([]bool, n)
	order := make([]int, 0, n)

	cur := startNode(g)
	visited[cur] = true
	order = append(order, cur)

	currentTimeMins := startTimeMins + g.Nodes[cur].DurationMin

	var totalDistM, totalTravelSec, totalServiceSec, totalWaitSec int
	totalServiceSec = g.Nodes[cur].DurationMin * 60

	for len(order) < n {
		next := -1
		nextDurSec := math.MaxInt

		// Проход 1: ближайший допустимый узел (укладывается в временное окно).
		for i := 0; i < n; i++ {
			if visited[i] {
				continue
			}
			e := g.Edges[cur][i]
			arrivalMins := currentTimeMins + e.DurationSec/60
			if feasible(g.Nodes[i], arrivalMins) && e.DurationSec < nextDurSec {
				nextDurSec = e.DurationSec
				next = i
			}
		}

		// Проход 2: запасной вариант — ближайший узел без учёта временного окна.
		if next == -1 {
			nextDurSec = math.MaxInt
			for i := 0; i < n; i++ {
				if visited[i] {
					continue
				}
				if e := g.Edges[cur][i]; e.DurationSec < nextDurSec {
					nextDurSec = e.DurationSec
					next = i
				}
			}
		}

		if next == -1 {
			break
		}

		e := g.Edges[cur][next]
		node := g.Nodes[next]

		arrivalMins := currentTimeMins + e.DurationSec/60
		waitMins := 0
		if node.WindowStart >= 0 && arrivalMins < node.WindowStart {
			waitMins = node.WindowStart - arrivalMins
		}

		totalDistM += e.DistanceM
		totalTravelSec += e.DurationSec
		totalServiceSec += node.DurationMin * 60
		totalWaitSec += waitMins * 60

		currentTimeMins = arrivalMins + waitMins + node.DurationMin
		visited[next] = true
		order = append(order, next)
		cur = next
	}

	return Result{
		Order:           order,
		TotalDistanceM:  totalDistM,
		TotalTravelSec:  totalTravelSec,
		TotalServiceSec: totalServiceSec,
		TotalWaitSec:    totalWaitSec,
	}, nil
}

// startNode выбирает начальный узел: с наиболее ранним открытием временного
// окна, или узел 0, если временные окна не заданы.
func startNode(g *Graph) int {
	best := 0
	for i, node := range g.Nodes {
		if node.WindowStart < 0 {
			continue
		}
		if g.Nodes[best].WindowStart < 0 ||
			node.WindowStart < g.Nodes[best].WindowStart {
			best = i
		}
	}
	return best
}

// feasible возвращает true, если обслуживание узла может начаться
// до закрытия его временного окна.
func feasible(node Node, arrivalMins int) bool {
	if node.WindowEnd < 0 {
		return true
	}
	return arrivalMins <= node.WindowEnd-node.DurationMin
}
  \end{minted}
\end{codelisting}

% ──────────────────────────────────────────────────────────────────────────────
% Листинг 2: HTTP-обработчики маршрутов (Go)
% ──────────────────────────────────────────────────────────────────────────────

\begin{codelisting}{HTTP-обработчики эндпоинтов маршрутов (handlers\_routes.go)}{app:lst_handlers_routes}
  \begin{minted}{go}
package http

import (
	"net/http"

	"github.com/gin-gonic/gin"
	routestory "planner-backend/internal/domain/route/story"
	"planner-backend/internal/platform/distance"
)

type RouteHandlers struct {
	story *routestory.Story
}

func NewRouteHandlers(st *routestory.Story) *RouteHandlers {
	return &RouteHandlers{story: st}
}

func (h *RouteHandlers) Create(c *gin.Context) {
	userID, ok := userIDFromCtx(c)
	if !ok {
		return
	}
	var req CreateRouteReq
	if err := c.ShouldBindJSON(&req); err != nil {
		c.JSON(http.StatusBadRequest, gin.H{"error": "bad request"})
		return
	}
	out, err := h.story.CreateDraft(
		c.Request.Context(), userID, req.Source, req.OrderedTaskIDs,
	)
	if err != nil {
		c.JSON(http.StatusBadRequest, gin.H{"error": err.Error()})
		return
	}
	c.JSON(http.StatusOK, out)
}

func (h *RouteHandlers) List(c *gin.Context) {
	userID, ok := userIDFromCtx(c)
	if !ok {
		return
	}
	out, err := h.story.List(c.Request.Context(), userID)
	if err != nil {
		c.JSON(http.StatusInternalServerError, gin.H{"error": "internal"})
		return
	}
	c.JSON(http.StatusOK, out)
}

// Optimize обрабатывает POST /api/routes/optimize.
// Получает матрицу расстояний от клиента, строит граф в памяти,
// запускает алгоритм оптимизации с учётом временных окон
// и сохраняет результирующий маршрут.
func (h *RouteHandlers) Optimize(c *gin.Context) {
	userID, ok := userIDFromCtx(c)
	if !ok {
		return
	}
	var req OptimizeRouteReq
	if err := c.ShouldBindJSON(&req); err != nil {
		c.JSON(http.StatusBadRequest, gin.H{"error": "bad request"})
		return
	}

	startTime := req.StartTimeMins
	if startTime <= 0 {
		startTime = 540 // по умолчанию: 09:00
	}

	// Преобразование клиентской матрицы расстояний в доменный тип.
	var matrix [][]distance.Edge
	if len(req.DistanceMatrix) > 0 {
		matrix = make([][]distance.Edge, len(req.DistanceMatrix))
		for i, row := range req.DistanceMatrix {
			matrix[i] = make([]distance.Edge, len(row))
			for j, cell := range row {
				matrix[i][j] = distance.Edge{
					DistanceM:   cell.DistanceM,
					DurationSec: cell.DurationSec,
				}
			}
		}
	}

	out, err := h.story.Optimize(
		c.Request.Context(), userID, req.TaskIDs, startTime, matrix,
	)
	if err != nil {
		c.JSON(http.StatusBadRequest, gin.H{"error": err.Error()})
		return
	}
	c.JSON(http.StatusOK, out)
}

func (h *RouteHandlers) Get(c *gin.Context) {
	userID, ok := userIDFromCtx(c)
	if !ok {
		return
	}
	routeID := c.Param("id")
	out, found, err := h.story.Get(c.Request.Context(), userID, routeID)
	if err != nil {
		c.JSON(http.StatusBadRequest, gin.H{"error": err.Error()})
		return
	}
	if !found {
		c.JSON(http.StatusNotFound, gin.H{"error": "not found"})
		return
	}
	c.JSON(http.StatusOK, out)
}

func (h *RouteHandlers) Delete(c *gin.Context) {
	userID, ok := userIDFromCtx(c)
	if !ok {
		return
	}
	routeID := c.Param("id")
	found, err := h.story.Delete(c.Request.Context(), userID, routeID)
	if err != nil {
		c.JSON(http.StatusInternalServerError, gin.H{"error": "internal"})
		return
	}
	if !found {
		c.JSON(http.StatusNotFound, gin.H{"error": "not found"})
		return
	}
	c.JSON(http.StatusOK, gin.H{"ok": true})
}

func (h *RouteHandlers) DeleteAll(c *gin.Context) {
	userID, ok := userIDFromCtx(c)
	if !ok {
		return
	}
	if err := h.story.DeleteAll(c.Request.Context(), userID); err != nil {
		c.JSON(http.StatusInternalServerError, gin.H{"error": "internal"})
		return
	}
	c.JSON(http.StatusOK, gin.H{"ok": true})
}

func (h *RouteHandlers) Rename(c *gin.Context) {
	userID, ok := userIDFromCtx(c)
	if !ok {
		return
	}
	routeID := c.Param("id")
	var req RenameRouteReq
	if err := c.ShouldBindJSON(&req); err != nil || req.Name == "" {
		c.JSON(http.StatusBadRequest, gin.H{"error": "name required"})
		return
	}
	found, err := h.story.Rename(
		c.Request.Context(), userID, routeID, req.Name,
	)
	if err != nil {
		c.JSON(http.StatusBadRequest, gin.H{"error": err.Error()})
		return
	}
	if !found {
		c.JSON(http.StatusNotFound, gin.H{"error": "not found"})
		return
	}
	c.JSON(http.StatusOK, gin.H{"ok": true})
}
  \end{minted}
\end{codelisting}

% ──────────────────────────────────────────────────────────────────────────────
% Листинг 3: API-клиент с автоматическим обновлением токена (TypeScript)
% ──────────────────────────────────────────────────────────────────────────────

\begin{codelisting}{Ключевые функции API-клиента: авторизованные запросы и обновление токена (client.ts)}{app:lst_api_client}
  \begin{minted}{typescript}
// Единственный глобальный промис обновления токена — предотвращает
// параллельные запросы /auth/refresh при одновременных 401-ошибках.
let refreshSessionPromise: Promise<AuthSession> | null = null;

// optimizeRoute передаёт идентификаторы задач и предварительно вычисленную
// матрицу расстояний Яндекс.Карт на бэкенд.
// Передача матрицы из buildYandexDistanceMatrix() обеспечивает согласованность
// данных оптимизации с отображением на карте.
export async function optimizeRoute(
  taskIds: string[],
  options: AuthorizedRequestOptions,
  startTimeMins = 540,
  distanceMatrix?: DistanceCell[][],
): Promise<SavedRouteDetails> {
  const route = await requestWithAuth<ApiRouteFull>(
    '/routes/optimize',
    {
      method: 'POST',
      body: JSON.stringify({ taskIds, startTimeMins, distanceMatrix }),
    },
    options,
  );
  return mapRouteDetails(route);
}

// requestWithAuth выполняет HTTP-запрос с токеном доступа.
// При получении 401 однократно обновляет токен и повторяет запрос.
async function requestWithAuth<T>(
  path: string,
  init: RequestInit,
  options: AuthorizedRequestOptions,
): Promise<T> {
  try {
    return await request<T>(path, init, options.session.accessToken);
  } catch (error) {
    if (!(error instanceof ApiError) || error.status !== 401) {
      throw error;
    }
  }
  try {
    const nextSession = await refreshSession(options);
    return request<T>(path, init, nextSession.accessToken);
  } catch (error) {
    await options.onUnauthorized();
    throw error;
  }
}

// refreshSession обновляет пару токенов через /auth/refresh.
// Дедупликация через единственный промис предотвращает гонку запросов.
async function refreshSession(options: AuthorizedRequestOptions) {
  if (!refreshSessionPromise) {
    refreshSessionPromise = request<TokenPairResponse>('/auth/refresh', {
      method: 'POST',
      body: JSON.stringify({ refreshToken: options.session.refreshToken }),
    })
      .then((pair) => {
        const nextSession = buildSession(
          pair,
          getEmailFromToken(pair.accessToken) ?? options.session.email,
        );
        options.onSessionChange(nextSession);
        return nextSession;
      })
      .finally(() => {
        refreshSessionPromise = null;
      });
  }
  return refreshSessionPromise;
}

// request — базовый HTTP-клиент: устанавливает заголовки Content-Type,
// Accept и Authorization, возвращает типизированный результат или
// выбрасывает ApiError с HTTP-статусом.
async function request<T = void>(
  path: string,
  init: RequestInit,
  accessToken?: string,
): Promise<T> {
  const headers = new Headers(init.headers);
  if (!headers.has('Content-Type') && init.body) {
    headers.set('Content-Type', 'application/json');
  }
  if (!headers.has('Accept')) {
    headers.set('Accept', 'application/json');
  }
  if (accessToken) {
    headers.set('Authorization', `Bearer ${accessToken}`);
  }
  const response = await fetch(`${API_BASE_URL}${path}`, { ...init, headers });
  if (!response.ok) {
    const payload = await parseErrorPayload(response);
    throw new ApiError(payload?.error ?? 'Request failed', response.status);
  }
  if (response.status === 204) return undefined as T;
  const text = await response.text();
  if (!text) return undefined as T;
  return JSON.parse(text) as T;
}
  \end{minted}
\end{codelisting}

% ──────────────────────────────────────────────────────────────────────────────
% Листинг 4: построение матрицы расстояний через Яндекс Карты (TypeScript)
% ──────────────────────────────────────────────────────────────────────────────

\begin{codelisting}{Построение матрицы расстояний через Яндекс Маршруты JS API (routeOptimizer.ts)}{app:lst_route_optimizer}
  \begin{minted}{typescript}
import { Task } from '../types/task';
import { loadYandexMaps } from './yandexMaps';

export interface DistanceCell {
  distanceM: number;
  durationSec: number;
}

/**
 * Строит матрицу расстояний n×n для списка задач через Яндекс Маршруты.
 * Используется тот же движок, что и для визуализации маршрута на карте,
 * поэтому данные оптимизации согласованы с отображением.
 *
 * matrix[i][j] = стоимость проезда от tasks[i] до tasks[j].
 * Диагональные элементы (i === j) равны нулю.
 */
export async function buildYandexDistanceMatrix(
  tasks: Task[],
  routingMode: 'auto' | 'pedestrian' | 'masstransit' = 'auto',
): Promise<DistanceCell[][]> {
  const n = tasks.length;
  if (n === 0) return [];

  await loadYandexMaps();

  // masstransit через multiRouter не возвращает надёжные данные по времени —
  // используем 'auto' как приближение.
  const mode = routingMode === 'masstransit' ? 'auto' : routingMode;

  const fallbackSec = 99_999;
  const matrix: DistanceCell[][] = Array.from({ length: n }, (_, i) =>
    Array.from({ length: n }, (__, j) =>
      i === j
        ? { distanceM: 0, durationSec: 0 }
        : { distanceM: 0, durationSec: fallbackSec },
    ),
  );

  // Запрашиваем все пары параллельно: n*(n-1) запросов.
  const pairs: [number, number][] = [];
  for (let i = 0; i < n; i++)
    for (let j = 0; j < n; j++)
      if (i !== j) pairs.push([i, j]);

  await Promise.all(
    pairs.map(async ([i, j]) => {
      const from = tasks[i];
      const to = tasks[j];
      if (from.latitude == null || to.latitude == null) return;

      try {
        const { distanceM, durationSec } = await new Promise<{
          distanceM: number;
          durationSec: number;
        }>((resolve, reject) => {
          const mr = new (window.ymaps as any).multiRouter.MultiRoute(
            {
              referencePoints: [
                [from.latitude!, from.longitude!],
                [to.latitude!, to.longitude!],
              ],
              params: { routingMode: mode },
            },
            {},
          );

          mr.model.events.once('requestsuccess', () => {
            try {
              const activeRoute: any = mr.getActiveRoute();
              if (activeRoute) {
                const distObj: any = activeRoute.properties.get('distance');
                const durObj: any = activeRoute.properties.get('duration');
                resolve({
                  distanceM: Math.round(
                    typeof distObj?.value === 'number' ? distObj.value : 0,
                  ),
                  durationSec: Math.round(
                    typeof durObj?.value === 'number'
                      ? durObj.value
                      : fallbackSec,
                  ),
                });
                return;
              }

              // Запасной путь: данные из модельных маршрутов напрямую.
              const modelRoutes: any[] = mr.model.getRoutes();
              if (modelRoutes.length > 0) {
                const legs: any[] = modelRoutes[0].getLegs();
                let totalDistM = 0;
                let totalDurSec = 0;
                legs.forEach((leg: any) => {
                  const d: any = leg.properties?.get?.('distance');
                  const dur: any = leg.properties?.get?.('duration');
                  totalDistM += typeof d?.value === 'number' ? d.value : 0;
                  totalDurSec +=
                    typeof dur?.value === 'number' ? dur.value : 0;
                });
                resolve({
                  distanceM: Math.round(totalDistM),
                  durationSec:
                    totalDurSec > 0
                      ? Math.round(totalDurSec)
                      : fallbackSec,
                });
                return;
              }
              reject(new Error('no route data'));
            } catch (e) {
              reject(e);
            }
          });

          mr.model.events.once('requestfail', () => {
            reject(new Error('routing failed'));
          });
        });

        matrix[i][j] = { distanceM, durationSec };
      } catch {
        // Оставляем fallback: алгоритм избегает недостижимых пар.
      }
    }),
  );

  return matrix;
}
  \end{minted}
\end{codelisting}
